% Source pp. 80-92 (Lectures7-12.pdf).
\section[Effective Action. Loop expansion]{Effective Action. Loop expansion\protect\footnote{PS Sect.~11.3--11.5.}}

The generating functional for the proper vertices
\begin{equation}\label{eq:Gammadef}
\Gamma[\phi]=\sum_{n=2}^\infty\frac1{n!}\int \dd^4x_1\cdots \dd^4x_n\,\Gamma^{(n)}(x_1,\dots,x_n)\,\phi(x_1)\cdots\phi(x_n)
\end{equation}
is called the \emph{Effective Action}. Here $-\Gamma^{(n)}(x_1,\dots,x_n)$ are
defined as the sums of one-particle irreducible diagrams for $n$-point
correlation functions, with amputated legs. It turns out that $\Gamma[\phi]$ is
related to the functional $W[J]$ by the functional version of the
\emph{Legendre transform}.

The Legendre transform must be familiar from classical mechanics and classical
thermodynamics. In the latter context it relates different thermodynamic
functions. For example, consider a magnet which in the absence of an external
field is described by a Hamiltonian $\Ham$, and let $\mu$ be its magnetic
moment. If we add an external field $B$, its internal energy is
\[
U(B)=\langle\Ham\rangle-MB
\]
where $M:=\langle\mu\rangle$ is the magnetization,
\[
M=-\frac{\partial F}{\partial B}
\]
and $F=U-TS$ is the free energy. One can solve the above equation for $B$:
$B=B(M)$, and define the thermodynamic potential (Gibbs free energy)
\[
\Phi(M)=\Bigl[F(B)+MB\Bigr]_{B=B(M)}=\langle\Ham\rangle_{B(M)}-TS\,,
\qquad\text{so that}\qquad
B(M)=\frac{\partial\Phi(M)}{\partial M}\,.
\]
In particular the equation
\[
\frac{\partial\Phi}{\partial M}=0
\]
determines spontaneous magnetization in the absence of external field. The
transformation relating $\Phi(M)$ to $F(B)$ is the \emph{Legendre transform}.

Note that the above functional $Z[J]$ is interpreted as the partition function
of a field statistical system with the Hamiltonian
\[
\mathcal A[\varphi]-\int J(x)\varphi(x)\dd^4x\,,
\]
i.e.\ the source $J(x)$ plays the role of the external field in the above
example. The generating functional $W[J]=\log\zeta[J]$ is then interpreted as
minus free energy,
\[
W[J]=-F[J]\,.
\]
In particular
\begin{equation}\label{eq:phiJ}
\phi(x,[J])\equiv\bigl\langle\varphi(x)\bigr\rangle_J=\frac{\delta W[J]}{\delta J(x)}\,,
\end{equation}
where $\langle\cdots\rangle_J$ denotes the expectation value in the presence
of the source $J$. Let us assume that the above equation can be solved for $J$,
defining the inverse functional
\[
J=J(x,[\phi]):\qquad J\bigl(x,[\phi(x,[J_0(x)])]\bigr)=J_0(x)\,,\qquad
\phi\bigl(x,[J(x,[\phi_0])]\bigr)=\phi_0(x)\,,
\]
for any functions $J_0(x)$ or $\phi_0(x)$. One can define then the analog of
the thermodynamic potential
\[
\Gamma[\phi]=\left[-W[J]+\int J(x)\phi(x)\dd^4x\right]_{J(x)=J(x,[\phi])}
\]
It is not difficult to show that for $\Gamma[\phi]$ thus defined the relation
holds
\begin{equation}\label{eq:JGamma}
J(x,[\phi])=\frac{\delta\Gamma[\phi]}{\delta\phi(x)}\,.
\end{equation}
Indeed, taking the variation of
\[
\Gamma[\phi]+W[J]-\int J(x)\phi(x)\dd^4x=0
\]
(where it is understood that $\phi$ and $J$ are related through
\eqref{eq:phiJ}) we have
\[
\int \dd^4x\,\left[\frac{\delta\Gamma}{\delta\phi(x)}\,\delta\phi(x)+\frac{\delta W}{\delta J(x)}\,\delta J(x)-\phi(x)\,\delta J(x)-J(x)\,\delta\phi(x)\right]=0
\]
The second and the third terms here cancel each other in view of
\eqref{eq:phiJ} and we obtain \eqref{eq:JGamma}.

It can be shown that the functional $\Gamma[\phi]$ coincides with the
generating functional of the proper vertices. Namely, its functional Taylor
expansion in $\phi$, Eq.~\eqref{eq:Gammadef}, has the proper vertices
$\Gamma^{(n)}$ as its coefficient functions. Let us see how it works for
lower-order correlation functions, using $\varphi^4$ theory as an example. We
have
\begin{multline*}
W[J]=\frac12\int \dd^4x_1\dd^4x_2\,W(x_1-x_2)\,J(x_1)J(x_2)\\
+\frac1{4!}\int\left[\prod_{i=1}^4\dd^4x_i\right]\,W^{(4)}(x_1,\dots,x_4)\,J(x_1)J(x_2)J(x_3)J(x_4)+O(J^6)
\end{multline*}
and
\begin{multline*}
\Gamma[\phi]=\frac12\int \dd^4x_1\dd^4x_2\,\Gamma^{(2)}(x_1-x_2)\,\phi(x_1)\phi(x_2)\\
+\frac1{4!}\int\left[\prod_{i=1}^4\dd^4x_i\right]\,\Gamma^{(4)}(x_1,\dots,x_4)\,\phi(x_1)\phi(x_2)\phi(x_3)\phi(x_4)+O(\phi^6)\,.
\end{multline*}
Using the equation $\phi(x,[J])=\delta W/\delta J(x)$ we can express $\phi$ in
terms of $J$ as
\begin{multline*}
\phi(x,[J])=\int \dd^4x_1\,W(x-x_1)\,J(x_1)\\
+\frac1{3!}\int\left[\prod_{i=1}^3\dd^4x_i\right]\,W^{(4)}(x,x_1,x_2,x_3)\,J(x_1)J(x_2)J(x_3)+O(J^5)\,.
\end{multline*}
Similarly, using $J(x,[\phi])=\delta\Gamma/\delta\phi(x)$, $J$ is expressed
through $\phi$ as
\begin{multline*}
J(x,[\phi])=\int \dd^4x_1\,\Gamma^{(2)}(x-x_1)\,\phi(x_1)\\
+\frac1{3!}\int\left[\prod_{i=1}^3\dd^4x_i\right]\,\Gamma^{(4)}(x,x_1,x_2,x_3)\,\phi(x_1)\phi(x_2)\phi(x_3)+O(\phi^5)\,.
\end{multline*}
These two relations must be inverse one to another. Substituting the first into
the second we have, schematically
\begin{multline*}
J(x)=\int_{y,y'}\Gamma^{(2)}(x-y)\,W(y-y')\,J(y')\\
+\frac1{3!}\int_{y,y_1,y_2,y_3}\Gamma^{(2)}(x-y)\,W^{(4)}(y,y_1,y_2,y_3)\,J(y_1)J(y_2)J(y_3)\\
+\frac1{3!}\int_{\substack{y'_1,y'_2,y'_3\\y_1,y_2,y_3}}\Gamma^{(4)}(x,y'_1,y'_2,y'_3)\,W(y'_1-y_1)\,W(y'_2-y_2)\\
\times W(y'_3-y_3)\,J(y_1)J(y_2)J(y_3)+O(J^5)\,.
\end{multline*}
As this must be an identity, we find again
\[
\Gamma^{(2)}(x-x')=S(x-x')\,,
\]
and
\[
W^{(4)}(x_1,\dots,x_4)=-\int_{y_1,\dots,y_4}\Gamma^{(4)}(y_1,\dots,y_4)\,W(y_1-x_1)\,W(y_2-x_2)\,W(y_3-x_3)\,W(y_4-x_4)\,,
\]
which coincide with our definitions of the proper vertices $\Gamma^{(2)}$ and
$\Gamma^{(4)}$.

\begin{exercise}
Carry out the Legendre transform to the order 6 in $J$ and $\phi$. Find
$W^{(6)}$ in terms of $\Gamma^{(6)}$, $\Gamma^{(4)}$, $\Gamma^{(2)}$.
\end{exercise}

\subsection*{Semiclassical (or Loop) Expansion}

There is a useful way to reorganize the diagrams contributing to $W[J]$ or
$\Gamma[\phi]$ known as the loop expansion. Generally, classical limit appears
as the leading term if one tries to estimate the functional integral using the
saddle point approximation. Then, the integral is dominated by the field
configurations close to the stationary point determined by solving the
equation $\delta\mathcal A[\varphi]=0$, i.e.\ the equations of motion of
classical field theory defined by the action $\mathcal A[\varphi]$. With this
in mind, let us restore the quantum parameter $\hbar$ (previously set to 1) in
the definition of the generating functional
\begin{equation}\label{eq:ZJhbar}
Z[J]=\int\bigl[D\varphi\bigr]\,e^{-(\mathcal A[\varphi]-\int_xJ(x)\varphi(x))/\hbar}
\end{equation}
The semiclassical approximation becomes exact in the limit $\hbar\to0$, while
corrections to the leading semiclassical contribution can be developed in the
form of a power series in $\hbar$. Note that as in statistical mechanics
interpretation $\hbar\sim kT$, and we can also talk here about the
low-temperature expansion. For this reason it is convenient to retain $\hbar$
as the formal parameter.\footnote{It is of course always possible to choose
units in which the actual Planck's constant is 1. Then formal parameter named
$\hbar$ in \eqref{eq:ZJhbar} is dimensionless, and it looks like an additional
parameter of the theory. In fact, such parameter can always be traded for
coupling constants determining the strength of interaction. For example, the
$\varphi^4$ action can be written as
$\mathcal A=\frac1\lambda\int\bigl(\frac12(\partial\chi)^2+\frac{m^2}2\chi^2+\frac1{4!}\chi^4\bigr)\dd^4x$
in terms of the rescaled field $\chi=\sqrt\lambda\,\varphi$.}

For the purpose of the present analysis I will consider a generic action of the
form
\begin{equation}\label{eq:genericA}
\mathcal A=\int\left[\frac12\,(\partial\varphi)^2+\frac{m^2}2\,\varphi^2+V(\varphi)\right]\,,
\end{equation}
with $V(\varphi)$ of general form
\[
V(\varphi)=\frac{\lambda_3}{3!}\,\varphi^3+\frac{\lambda_4}{4!}\,\varphi^4+\cdots
\]
with constants $\lambda_k$, although I will often give illustrations in terms
of $\varphi^4$ theory.

\subsubsection*{The leading (tree) approximation}

For small $\hbar$ we calculate the above integral \eqref{eq:ZJhbar} for $Z[J]$
using a saddle-point approximation. That is, one first finds the minimum of the
action functional
\[
\mathcal A_J[\varphi]=\mathcal A[\varphi]-\int_xJ(x)\varphi(x)\,,
\]
i.e.\ the solution of the classical equation of motion
\[
\frac{\delta\mathcal A[\varphi]}{\delta\varphi(x)}=J(x)\,.
\]
The solution depends on $J$; I will denote it as
\[
\varphi_c(x)=\varphi_c(x,[J])\,.
\]
Then in the leading approximation
\[
\log Z[J]=\frac1\hbar\Bigl(-\mathcal A[\varphi_c]+\int_xJ(x)\varphi_c(x)\Bigr)+O(\hbar^0)
\]
When $\hbar$ is explicit it is convenient to normalize the generating
functional $W[J]$ so that
\[
\frac{Z[J]}{Z[0]}=e^{W[J]/\hbar}\,.
\]
We write
\[
W[J]=W_0[J]+\hbar\,W_1[J]+\cdots\,.
\]
In this notation we have
\begin{equation}\label{eq:W0}
W_0[J]=-\mathcal A[\varphi_c]+\int J(x)\varphi_c(x)\,,
\end{equation}
where it is understood that $\varphi_c=\varphi_c[J]$.

The relation of this leading semiclassical approximation to the ordinary
perturbation theory (expansion in $\lambda_i$) is recovered by expanding the
solution $\varphi_c$ in the powers of $J$. In view of the chosen form of the
action \eqref{eq:genericA} the classical equation has the form
\[
\bigl(m^2-\Delta_x\bigr)\varphi_c(x)+V'\bigl(\varphi_c(x)\bigr)=J(x)\,.
\]
This can be transformed to the integral equation
\[
\varphi_c(x)=\int_{x'}D(x-x')\,J(x')-\int_{x'}D(x-x')\,V'\bigl(\varphi_c(x')\bigr)
\]
which can be solved iteratively, with the last term with $V'$ taken as the
perturbation. For example, if
\[
V=\frac{\lambda}{4!}\,\varphi^4\,,\qquad V'=\frac{\lambda}{3!}\,\varphi^3\,,
\]
we find
\begin{multline*}
\varphi_c(x,[J])=\int_{x'}D(x-x')\,J(x')\\
-\frac{\lambda}{3!}\int_{x',x_1,x_2,x_3}D(x-x')\,D(x'-x_1)\,D(x'-x_2)\,D(x'-x_3)\,J(x_1)J(x_2)J(x_3)+\cdots\,.
\end{multline*}
Obviously, this generates the expansion in $J$. Also, this series is
represented diagrammatically as
\[
\varphi_c(x,[J])=
x\,\dia{\node[xvtx] at (0,0) {}; \draw[prop] (0,0) -- (3,0) node[vtx] {};}
+x\,\dia{\node[xvtx] at (0,0) {}; \draw[prop] (0,0) -- (2,0); \foreach \a in {50,0,-50} \draw[prop] (2,0) -- ++(\a:1.4) node[vtx] {}; \node[vtx] at (2,0) {};}
+x\,\dia{\node[xvtx] at (0,0) {}; \draw[prop] (0,0) -- (2,0); \foreach \a in {50,0,-50} \draw[prop] (2,0) -- ++(\a:1.4) node[vtx] {}; \node[vtx] at (2,0) {};
  \path (2,0) ++(50:1.4) coordinate (b); \foreach \a in {120,60,10} \draw[prop] (b) -- ++(\a:0.9) node[vtx] {}; \node[vtx] at (b) {};}
+\cdots\,.
\]
Here I mark the ``external point'' $x$ where $\varphi_c(x)$ sits by a cross to
distinguish it from the one-leg vertex
\[
\dia{\draw[prop] (-2,0) -- (0,0) node[vtx] {};}=\int \dd^4x\,J(x)\,.
\]
It is easy to see that the diagrams appearing here have no loops (i.e.\ one
cannot travel along the diagram and come back to the starting point never going
the same line twice). Such diagrams are called the \emph{tree diagrams}.

The easiest way to obtain the functional $W_0[J]$ \eqref{eq:W0} is to observe
that
\[
\frac{\delta W_0}{\delta J}=\left[-\frac{\delta\mathcal A}{\delta\varphi_c}+J\right]\frac{\delta\varphi_c}{\delta J}+\varphi_c=\varphi_c\,,
\]
where the fact that $\varphi_c$ solves the classical equation is used.
Integrating the above expansion of $\varphi_c(x,[J])$ term by term we find
\[
W_0[J]=\underbrace{\dia{\pic{line};}}_{W_0^{(2)}}
+\underbrace{\dia{\pic{legblob={0}{45,135,225,315}{1.4}{fill=none}}; \node[vtx] at (0,0) {};}}_{W_0^{(4)}}
+\underbrace{\Bigl(\dia{\draw[prop] (-0.8,0) -- (0.8,0); \foreach \a in {60,180,300} \draw[prop] (-0.8,0) -- ++(\a:1.1) node[vtx] {};
  \foreach \a in {120,0,240} \draw[prop] (0.8,0) -- ++(\a:1.1) node[vtx] {}; \node[vtx] at (-0.8,0) {}; \node[vtx] at (0.8,0) {};}
 +\dia{\pic{legblob={0}{0,60,120,180,240,300}{1.3}{fill=none}}; \node[vtx] at (0,0) {};}\Bigr)}_{W_0^{(6)}}+\cdots,
\]
i.e.\ the leading semiclassical approximation for $W$ is just the sum of all
connected tree diagrams. For this reason the leading approximation is also
called the \emph{tree approximation}.

Let us find the effective action in this approximation. By the definition of
the Legendre transform we have to set
\[
\phi=\frac{\delta W_0}{\delta J}=\varphi_c[J]\,.
\]
Then the effective action in this approximation is
\[
\Gamma_0[\phi]=-W_0[J]+J\phi=\mathcal A[\varphi_c]-J\varphi_c+J\varphi_c=\mathcal A[\phi]\,,
\]
i.e.\ in the leading semiclassical approximation the effective action
$\Gamma_0[\phi]$ coincides with the action $\mathcal A[\phi]$. In particular
\[
\Gamma^{(2)}_0(x,x')=S_0(x-x')=\bigl(m^2-\Delta_x\bigr)\,\delta(x-x')\,;
\]
and assuming for example that
$V(\varphi)=\frac{\lambda_4}{4!}\varphi^4+\frac{\lambda_6}{6!}\varphi^6+\cdots$,
\[
\Gamma^{(4)}_0=-\ \dia{\foreach \a in {45,135,225,315} \draw[prop,densely dashed] (0,0) -- (\a:1.4); \node[vtx,rectangle,minimum size=3.5pt] at (0,0) {};}\ ;\qquad
\Gamma^{(6)}_0=-\ \dia{\foreach \a in {0,60,120,180,240,300} \draw[prop,densely dashed] (0,0) -- (\a:1.4); \node[vtx,rectangle,minimum size=3.5pt] at (0,0) {};}\ ;\qquad\dots,
\]
where I have shown where the amputated legs used to be by the dashed lines.
Clearly, these are the simplest one-particle irreducible amputated diagrams for
$\Gamma^{(n)}$.

\subsubsection*{The first order (one-loop) correction}

Let us now turn to the first subleading terms $W_1[J]$ and $\Gamma_1[\phi]$ in
the expansions
\[
\frac{W[J]}\hbar=\frac{W_0[J]}\hbar+W_1[J]+O(\hbar)\,,\qquad
\frac{\Gamma[\phi]}\hbar=\frac{\Gamma_0[\phi]}\hbar+\Gamma_1[\phi]+O(\hbar)\,.
\]
As we have previously seen, $W_0[J]$ is given by the sum of all connected tree
diagrams, while $\Gamma_0[\phi]$ coincides with the action $\mathcal A[\phi]$.

To evaluate the first correction to the effective action one must take into
account the gaussian fluctuations around the classical solution $\varphi_c$. In
the functional integral \eqref{eq:ZJhbar} defining $Z[J]$ set
\[
\varphi(x)=\varphi_c(x)+\chi(x)
\]
and expand the action $\mathcal A_J$ to the terms quadratic in $\chi$,
\[
\mathcal A_J[\varphi]=\mathcal A[\varphi_c]-J\varphi_c+\frac{\delta\mathcal A}{\delta\varphi}\Big|_{\varphi=\varphi_c}\chi-J\chi
+\int_{x,x'}\frac12\,\frac{\delta^2\mathcal A}{\delta\varphi(x)\,\delta\varphi(x')}\Big|_{\varphi=\varphi_c}\chi(x)\chi(x')+O(\chi^3)\,.
\]
As usual, the terms linear in $\chi$ cancel in view of the classical equation
$\varphi_c$ satisfies, and we obtain in this approximation
\[
Z[J]=e^{W_0[J]/\hbar}\int\bigl[D\chi\bigr]\,e^{-\mathcal A_1[\chi]/\hbar}\,,
\]
where
\begin{equation}\label{eq:A1}
\mathcal A_1[\chi]=\int \dd^4x\,\left[\frac12\,(\partial\chi)^2+\frac{m^2}2\,\chi^2+\frac12\,V''(\varphi_c)\,\chi^2\right]
\end{equation}
Since the higher-order terms in $\chi$ are neglected the integral over $\chi$
is gaussian and it can be evaluated explicitly in terms of the determinant of
the operator
\[
K(x,x')=\frac{\delta^2\mathcal A}{\delta\varphi(x)\,\delta\varphi(x')}\Big|_{\varphi=\varphi_c}
=\bigl(m^2-\Delta_x\bigr)\,\delta(x-x')+V''(\varphi_c)\,\delta(x-x')\,.
\]

\begin{exercise}
Evaluate $n$-dimensional gaussian integral
\[
\prod_{i=1}^n\left[\int_{-\infty}^\infty \dd\xi_i\right]\,\exp\left\{-\frac12\sum_{i,j=1}^ns_{ij}\,\xi_i\xi_j\right\}
\]
where $s_{ij}$ are elements of a symmetric positive matrix, and show that it
equals
\[
(2\pi)^{n/2}\,\bigl[\det(s_{ij})\bigr]^{-1/2}\,.
\]
\end{exercise}

Actually, we are interested in the ratio $Z[J]/Z[0]$, therefore we need to
evaluate the ratio of the determinants
\begin{multline*}
\left[\frac{\det\bigl(S_0(x-x')+\delta(x-x')\,V''(x')\bigr)}{\det\bigl(S_0(x-x')\bigr)}\right]^{-1/2}
\\=\exp\left[-\frac12\log\det\Bigl(\delta(x-x')+D(x-x')\,V''\bigl(\varphi_c(x')\bigr)\Bigr)\right]
\end{multline*}
This can be transformed with the help of the identity, valid for any positive
Hermitian operator $M$:
\[
\log\bigl(\det M\bigr)=\tr\bigl(\log M\bigr)\,,
\]
which is established by bringing $M$ to the diagonal form,
\[
\log\Bigl(\prod_im_i\Bigr)=\sum_i\log m_i\,.
\]
Using this property, we obtain
\begin{equation}\label{eq:W1}
W_1[J]=-\frac12\,\tr\log\bigl(I+A(\varphi_c)\bigr)\,,
\end{equation}
where $A$ is an integral operator with the kernel
\[
A(x,x')=D(x-x')\,V''\bigl(\varphi_c(x')\bigr)\,,
\]
and $I$ is the identity operator, $I(x,x')=\delta(x-x')$. The logarithm expands
as
\[
\log\bigl(I+A\bigr)=A-\frac{A^2}2+\frac{A^3}3-\dots-\frac{(-A)^n}n+\cdots\,.
\]
When taking the trace we obtain
\[
\tr A=\int_xD(x-x)\,V''(x)\,,\qquad
\tr A^2=\int_{x,x'}D(x-x')\,V''(x')\,D(x'-x)\,V''(x)\,,\qquad\text{etc.}
\]
Here $V''(x)$ is the shorthand for $V''(\varphi_c(x))$. Obviously, this
expansion can be represented in terms of diagrams
\[
W_1[J]=-\frac12\,\tr\log\bigl(I+A\bigr)
=\dia{\draw[line width=1.2pt] (0,0) circle (1); \draw[prop,fill=white] (0:1) circle (0.25);}
+\dia{\draw[line width=1.2pt] (0,0) circle (1); \foreach \a in {0,180} \draw[prop,fill=white] (\a:1) circle (0.25);}
+\dia{\draw[line width=1.2pt] (0,0) circle (1); \foreach \a in {30,150,270} \draw[prop,fill=white] (\a:1) circle (0.25);}
+\cdots\,,
\]
with the line standing for $D$, as usual, but with the special two-leg vertex
\[
\dia{\draw[prop] (-1.5,0) -- (1.5,0); \draw[prop,fill=white] (0,0) circle (0.4);}
=-\int \dd^4x\,V''\bigl(\varphi_c(x)\bigr)\,.
\]
A diagram with $n$ such vertices comes with the factor (implicit in the above
series)
\[
\frac1{2n}=\frac12\times\frac1n\,,
\]
$1/2$ already being a factor in \eqref{eq:W1}, and $1/n$ coming from the
expansion coefficient of the log function; this factor can be interpreted as
the symmetry factor of the $n$-vertex diagram. Of course the above diagram
expansion can be obtained directly from the Feynman rules for the action
\eqref{eq:A1}.

To obtain $W_1[J]$, substitute
\[
\varphi_c=\varphi_c[J]=
\dia{\node[xvtx] at (0,0) {}; \draw[prop] (0,0) -- (3,0) node[vtx] {};}
+\dia{\node[xvtx] at (0,0) {}; \draw[prop] (0,0) -- (2,0); \foreach \a in {50,0,-50} \draw[prop] (2,0) -- ++(\a:1.4) node[vtx] {}; \node[vtx] at (2,0) {};}
+\cdots\,,
\]
which, as we have seen, is a sum of the tree diagrams. For $\varphi^4$ theory
\[
V''(\varphi_c)=\frac\lambda2\,\varphi_c^2
\]
and so
\[
\dia{\draw[prop] (-1.5,0) -- (1.5,0); \draw[prop,fill=white] (0,0) circle (0.4);}
=\dia{\draw[prop] (-1.5,0) -- (1.5,0); \draw[prop] (-0.8,1.4) node[vtx] {} -- (0,0) -- (0.8,1.4) node[vtx] {}; \node[vtx] at (0,0) {};}
+\dia{\draw[prop] (-1.5,0) -- (1.5,0); \draw[prop] (-0.8,1.4) node[vtx] {} -- (0,0) -- (0.8,1.4); \node[vtx] at (0,0) {};
  \foreach \a in {120,60,10} \draw[prop] (0.8,1.4) -- ++(\a:0.8) node[vtx] {}; \node[vtx] at (0.8,1.4) {};}
+\cdots\,.
\]
Therefore $W_1$ is the sum of the diagrams of the form
\[
\dia{\draw[line width=1.2pt] (0,0) circle (1); \foreach \a in {30,150,270} \draw[prop,fill=white] (\a:1) circle (0.25);}
\ \longrightarrow\
\dia{\draw[prop] (0,0) circle (1); \foreach \a in {30,150,270} \node[vtx] at (\a:1) {};
  \draw[prop] (30:1) -- ++(60:0.9) node[vtx] {}; \draw[prop] (30:1) -- ++(0:0.9) node[vtx] {};
  \draw[prop] (270:1) -- ++(240:0.9) node[vtx] {}; \draw[prop] (270:1) -- ++(300:0.9) node[vtx] {};
  \draw[prop] (150:1) -- ++(150:0.9) coordinate (c) node[vtx] {}; \draw[prop] (150:1) -- ++(90:0.9) node[vtx] {};
  \foreach \a in {200,150,100} \draw[prop] (c) -- ++(\a:0.7) node[vtx] {};}
\]
which have exactly one loop. This is why the $\sim\hbar$ correction to $W$ is
called the ``one-loop correction''.

It is not difficult to find the effective action in this approximation. We
have
\[
\Gamma[\phi]=-W_0[J]-\hbar\,W_1[J]+\int J\phi+O(\hbar^2)\,,
\]
where we have to substitute $J=J[\phi]$. Note that by definition of the
Legendre transform, the r.h.s.\ of this equation is stationary with respect to
variations of $J$ with $\phi$ fixed. Therefore the correction $\sim\hbar$ to
the relation between $J$ and $\phi$ will produce a change of the order of
$\hbar^2$ to the r.h.s.\ of the above equation. Hence at order $\hbar$ we still
can replace $J$ by the solution of the equation
\[
\phi(x)=\varphi_c(x,[J])\,.
\]
Then, again $-W_0+\int J\phi=\mathcal A[\phi]$, and
\[
\Gamma[\phi]=\mathcal A[\phi]+\hbar\,\Gamma_1[\phi]+O(\hbar^2)\,,
\]
where
\[
\Gamma_1[\phi]=\frac12\,\tr\log\bigl(I+A(\phi)\bigr)
=-\Bigl(\dia{\draw[line width=1.2pt] (0,0) circle (1); \node[otimesv] at (0:1) {};}
+\dia{\draw[line width=1.2pt] (0,0) circle (1); \foreach \a in {0,180} \node[otimesv] at (\a:1) {};}
+\dia{\draw[line width=1.2pt] (0,0) circle (1); \foreach \a in {30,150,270} \node[otimesv] at (\a:1) {};}
+\cdots\Bigr)\,,
\]
with
\[
\dia{\node[otimesv] at (0,0) {};}
=-V''(\phi)\,.
\]
Note that all these diagrams are one-particle irreducible.

Similar analysis can be carried out in higher orders in $\hbar$. One finds that
at the order $\hbar^n$ only the diagrams with exactly $n$ loops contribute to
$W_n[J]$.

\subsection*{Effective Action vs Generating Functional}

We can now clarify the nature of the relation between $W[J]$ and
$\Gamma[\phi]$. Write again
\[
\Gamma[\phi]=\sum_{n=2}^\infty\frac1{n!}\int\Gamma^{(n)}(x_1,\dots,x_n)\,\phi(x_1)\cdots\phi(x_n)\,,
\]
and represent $\Gamma^{(n)}$ as an $n$-point blob
\[
\Gamma^{(n)}=\diac{\draw[prop,pattern=north east lines] (0,0) circle (1);
  \foreach \a/\l/\p in {180/x_1/left,125/x_2/above left,55/x_3/above right,0/x_4/right} {\node[vtx,minimum size=4pt] at (\a:1) {}; \node[\p] at (\a:1) {\small$\l$};}
  \draw[densely dotted] (200:1.35) arc[start angle=200,end angle=340,radius=1.35];}\,.
\]
The generating functional $W[J]$ is related to $\Gamma[\phi]$ by a functional
Legendre transform.
\[
W[J]=\Bigl[-\Gamma[\phi]+\int_xJ(x)\phi(x)\Bigr]_{\phi=\phi[J]}
\]
One first solves the equation
\[
J(x)=\frac{\delta\Gamma[\phi]}{\delta\phi(x)}
\]
with respect to $\phi(x)$. Exactly as was done when solving the classical
equations of motion, the solution is represented as a sum of tree diagrams
\[
\phi(x,[J])=
\dia{\node[xvtx] at (0,0) {}; \draw[prop] (0,0) -- (3.2,0) node[vtx] {}; \draw[prop,fill=white] (1.6,0) circle (0.45);}
+\dia{\node[xvtx] at (-2.3,0) {}; \draw[prop] (-2.3,0) -- (-0.7,0); \draw[prop,fill=white] (-1.5,0) circle (0.28);
  \pic{dressed={0.7}{60,0,-60}};}
+\cdots\,,
\]
this time with $W(x-x')$ (the inverse of $\Gamma^{(2)}$, i.e.\ the full
two-point function) standing instead of the propagators $D$, and the functions
$\Gamma^{(n)}$ playing the role of the vertices. In drawing the above diagrams
I have assumed that the interaction was even, $V(\varphi)=V(-\varphi)$, so that
odd order vertices are absent; of course, this limitation is not significant.
Again, integrating the equation
\[
\frac{\delta W[J]}{\delta J(x)}=\phi(x,[J])
\]
term by term in $J$, one finds that $W[J]$ is a sum of all tree diagrams, again
with $W$ for propagators and $\Gamma^{(n)}$ for vertices,
\[
W[J]=\dia{\draw[prop] (-1.6,0) node[vtx]{} -- (1.6,0) node[vtx]{}; \draw[prop,fill=white] (0,0) circle (0.45);}
+\diac{\pic{dressed={0.7}{45,135,225,315}};}
+\ \text{all tree diagrams}
\]
This is exactly what we expect to be the relation between the sum of all
connected diagrams and the sum of all one-particle irreducible diagrams: one
takes one-particle irreducible diagrams and connects them in a tree-like manner
to produce all connected diagrams.

\begin{hwproblem}{3}\label{hw:3}
In the $\varphi^3$ theory
\[
\mathcal A_I=\frac{\lambda_3}{3!}\int\varphi^3(x)\dd^4x
\]
carry out the Legendre transform of $W[J]$ to the order 4 in $J$ and $\phi$.
Define
\[
\hat\Gamma[\phi]:=\Gamma[\bar\phi+\phi]-\Gamma[\bar\phi]
\qquad\text{where}\qquad \bar\phi=W^{(1)}(x)
\]
and $\Gamma[\phi]$ is the Legendre transform of $W[J]$. Express $W^{(3)}$ and
$W^{(4)}$ in terms of $\hat\Gamma^{(4)}$, $\hat\Gamma^{(3)}$, $\hat\Gamma^{(2)}$.
Give diagrammatic representation of your result.
\end{hwproblem}
